% =========================================================================
% SEÇÃO 3: FUNDAMENTAÇÃO TEÓRICA
% =========================================================================
\section{Fundamentação Teórica}\label{sec:fundamentacao}

Esta seção formaliza os modelos matemáticos, teoremas estruturais e mecanismos algorítmicos que regem os problemas de Fluxo Máximo e Fluxo de Custo Mínimo.

\subsection{Base Matemática dos Grafos de Fluxo}

Uma rede de fluxo é representada por um digrafo $D = (V, A)$, onde $V$ denota o conjunto de vértices e $A$ o conjunto de arcos direcionados (veja~\cite{ahuja1993network, bondy2008graph}). A rede possui dois vértices terminais distinguidos: uma fonte $s \in V$ e um sumidouro $t \in V$. Cada arco $(u, v) \in A$ possui uma capacidade não negativa $c(u, v) \ge 0$, estipulada pela função $c: A \to \mathbb{R}_{\ge 0}$, que estabelece o limite superior de tráfego admissível no arco.

Um fluxo na rede $D$ é uma função $f: A \to \mathbb{R}_{\ge 0}$ que associa a cada arco $(u, v)$ um escalar $f(u, v)$. O fluxo $f$ é dito viável se satisfaz duas condições axiomáticas (veja~\cite{ford1956maximal, cook1998combinatorial}):
\begin{enumerate}
    \item \textbf{Restrição de Capacidade:} O fluxo em cada conexão é não negativo e não excede sua capacidade:
    \begin{equation}
        0 \le f(u, v) \le c(u, v), \quad \forall (u, v) \in A
    \end{equation}
    \item \textbf{Conservação de Fluxo:} Para todo vértice intermediário $v \in V \setminus \{s, t\}$, a soma dos fluxos que entram em $v$ é igual à soma dos fluxos que dele saem:
    \begin{equation}
        \sum_{\{u \mid (u, v) \in A\}} f(u, v) = \sum_{\{w \mid (v, w) \in A\}} f(v, w)
    \end{equation}
\end{enumerate}

O valor total do fluxo, denotado por $|f|$, quantifica a taxa líquida que emerge da fonte $s$:
\begin{equation}
    |f| = \sum_{\{v \mid (s, v) \in A\}} f(s, v) - \sum_{\{u \mid (u, s) \in A\}} f(u, s)
\end{equation}
Pela conservação nodal, esse valor coincide rigorosamente com o fluxo líquido absorvido pelo sumidouro $t$. O Problema do Fluxo Máximo consiste em determinar uma função viável $f$ que maximize $|f|$.

\subsection{Redes Residuais, Caminhos Aumentantes e Dualidade}

Dado um fluxo viável $f$, a rede residual $D_f = (V, A_f)$ quantifica a capacidade ociosa dos arcos originais e a possibilidade de cancelamento de fluxos previamente despachados (veja~\cite{ahuja1993network}). A capacidade residual $c_f(u, v)$ é definida por:
\begin{equation}
    c_f(u, v) = c(u, v) - f(u, v) + f(v, u)
\end{equation}
O conjunto $A_f$ é formado por todos os arcos direcionados $(u, v)$ tais que $c_f(u, v) > 0$. Um $st$-caminho aumentante é um caminho direcionado simples de $s$ para $t$ em $D_f$. Se existir tal caminho $P$, a capacidade residual gargalo $\delta = \min_{(u, v) \in P} c_f(u, v) > 0$ pode ser enviada ao longo de $P$, gerando um novo fluxo viável $f'$ com valor $|f'| = |f| + \delta$.

Um $st$-corte $(S, T)$ é uma partição de $V$ tal que $s \in S$ e $t \in T = V \setminus S$. A capacidade do corte é a soma das capacidades dos arcos direcionados de $S$ para $T$:
\begin{equation}
    c(S, T) = \sum_{u \in S, v \in T, (u, v) \in A} c(u, v)
\end{equation}

\begin{prop}[Teorema do Fluxo Máximo e Corte Mínimo --- Ford-Fulkerson, 1956]\label{prop:mfmc_rel}
Em qualquer rede de fluxo, o valor máximo de um $st$-fluxo viável é igual à capacidade mínima de um $st$-corte:
\begin{equation}
    \max_{f \text{ viável}} |f| = \min_{(S, T) \text{ corte}} c(S, T)
\end{equation}
Além disso, as seguintes condições são equivalentes: (i) $f$ é um fluxo máximo; (ii) $D_f$ não contém nenhum $st$-caminho aumentante; (iii) existe um corte $(S^*, T^*)$ tal que $|f| = c(S^*, T^*)$.
\end{prop}

\noindent \textbf{Demonstração Construtiva:} Para qualquer fluxo viável $f$ e qualquer corte $(S, T)$, tem-se $|f| \le c(S, T)$. Se $D_f$ não possui caminho aumentante de $s$ a $t$, seja $S^*$ o conjunto de vértices alcançáveis a partir de $s$ em $D_f$ e $T^* = V \setminus S^*$. Como não há caminho de $s$ a $t$, $s \in S^*$ e $t \in T^*$, de modo que $(S^*, T^*)$ é um corte válido. Para todo $(u, v) \in A$ com $u \in S^*$ e $v \in T^*$, deve-se ter $c_f(u, v) = 0$, implicando $f(u, v) = c(u, v)$. De forma análoga, se $u \in T^*$ e $v \in S^*$, tem-se $f(u, v) = 0$ (caso contrário, existiria capacidade residual reversa $c_f(v, u) > 0$, contradizendo a inalcançabilidade de $u$). O fluxo líquido através do corte satisfaz $|f| = c(S^*, T^*)$, certificando a otimalidade de $f$ e a minimalidade de $(S^*, T^*)$. \hfill $\blacksquare$

\begin{figure}[!ht]
\centering
\begin{tikzpicture}[thick]
    \draw[cut S] (-0.8, -1.8) rectangle (2.8, 1.8);
    \node[cut label S] at (-0.6, 1.65) {$\mathbf{S}$};

    \draw[cut T] (5.4, -1.8) rectangle (9.0, 1.8);
    \node[cut label T] at (8.8, 1.65) {$\mathbf{T}$};

    \node[source node] (s) at (0.0, 0.0) {$s$};
    \node[s node internal] (a) at (2.0, 0.8) {$a$};
    \node[s node internal] (b) at (2.0, -0.8) {$b$};

    \node[t node internal] (c) at (6.2, 0.8) {$c$};
    \node[t node internal] (d) at (6.2, -0.8) {$d$};
    \node[sink node] (t) at (8.2, 0.0) {$t$};

    \draw[line width=1.5pt, dashed, cutviolet] (4.1, -1.85) -- (4.1, 1.85)
        node[above=2.5pt, solid, font=\scriptsize\bfseries, text=cutviolet!90!black, fill=white, inner sep=2.5pt, rounded corners=3pt, draw=cutviolet!70] {Corte Mínimo $(S,T)$};

    \draw[sat edge] (a) -- node[edge lbl above, text=cutviolet!90!black] {$f=c$} (c);
    \draw[sat edge] (b) -- node[edge lbl below, text=cutviolet!90!black] {$f=c$} (d);

    \draw[flow edge] (s) -- (a);
    \draw[flow edge] (s) -- (b);
    \draw[flow edge] (c) -- (t);
    \draw[flow edge] (d) -- (t);
    \draw[flow edge] (a) -- (b);
    \draw[flow edge] (c) -- (d);
\end{tikzpicture}

\caption{Estrutura do corte mínimo $(S,T)$: arcos que cruzam de $S$ para $T$ (em violeta) estão saturados ($f = c$), e fluxos no sentido reverso de $T$ para $S$ são nulos.}
\label{fig:corte_minimo}
\end{figure}

\begin{prop}[Teorema da Integralidade]\label{prop:integralidade_rel}
Se a capacidade $c(a)$ for inteira para todo $a \in A$, então existe um fluxo máximo $f^*$ que é puramente inteiro, isto é, $f^*(a) \in \mathbb{Z}_{\ge 0}$ para todo arco $a \in A$.
\end{prop}

A preservação da integralidade sob capacidades inteiras é garantida pelo fato de que todo aumento ao longo de caminhos residuais em redes inteiras transporta um gargalo $\delta \in \mathbb{Z}_{> 0}$, viabilizando a modelagem exata de problemas combinatórios discretos.

\subsection{Algoritmos de Caminhos Aumentantes}

O algoritmo genérico de Ford-Fulkerson busca caminhos aumentantes via busca em profundidade (DFS). Sob capacidades inteiras, sua complexidade assintótica é $\mathcal{O}(|A| \cdot |f^*|)$, que depende do valor numérico das capacidades (pseudopolinomial). Para obter complexidade polinomial estrita, a literatura introduziu políticas estruturadas de seleção de caminhos.

A Figura~\ref{fig:ford_fulkerson_patologico} ilustra a patologia clássica do método ingênuo com DFS: capacidades de grande magnitude ($M = 10^6$) nos arcos externos e uma conexão transversal unitária ($c=1$) provocam alternâncias sucessivas de aumento unitário ($\delta = 1$), exigindo $2 \times 10^6$ passos. Em contrapartida, algoritmos guiados por distâncias mínimas em arcos (BFS) priorizam caminhos diretos de comprimento 2 ($s \to u \to t$ e $s \to v \to t$), saturando a rede em 2 passos (Edmonds-Karp) ou em uma única fase de fluxo bloqueador (Dinic), desconsiderando o arco transversal por conectar nós no mesmo nível ($d(u) = d(v) = 1$).

\begin{figure}[!ht]
\centering
\resizebox{0.92\textwidth}{!}{%
\begin{tikzpicture}[thick]
\begin{scope}[xshift=0cm]
  \tikzsubcaption{(2.5,2.6)}{(a) Rede patológica ($M = 10^6$)}
  \node[source node] (s) at (0, 1) {$s$};
  \node[flow node] (u) at (2.5, 2) {$u$};
  \node[flow node] (v) at (2.5, 0) {$v$};
  \node[sink node] (t) at (5, 1) {$t$};

  \draw[flow edge] (s) -- node[edge lbl above] {$M$} (u);
  \draw[flow edge] (s) -- node[edge lbl below] {$M$} (v);
  \draw[sat edge] (u) -- node[edge lbl right, text=red!80!black] {$c=1$} (v);
  \draw[flow edge] (u) -- node[edge lbl above] {$M$} (t);
  \draw[flow edge] (v) -- node[edge lbl below] {$M$} (t);
\end{scope}

\begin{scope}[xshift=7.2cm]
  \tikzsubcaption{(2.5,2.6)}{(b) Degeneração da DFS: $2M$ iterações}
  \node[source node] (s2) at (0, 1) {$s$};
  \node[flow node] (u2) at (2.5, 2) {$u$};
  \node[flow node] (v2) at (2.5, 0) {$v$};
  \node[sink node] (t2) at (5, 1) {$t$};

  \draw[tree edge, line width=1.8pt] (s2) -- node[edge lbl above] {$+1$} (u2);
  \draw[sat edge, line width=1.8pt] (u2) -- node[edge lbl right] {$+1$} (v2);
  \draw[tree edge, line width=1.8pt] (v2) -- node[edge lbl below] {$+1$} (t2);

  \draw[flow edge, opacity=0.3] (s2) -- (v2);
  \draw[flow edge, opacity=0.3] (u2) -- (t2);

  \node[font=\scriptsize, align=center, below=3pt] at (2.5,-0.5) {DFS alterna $s \to u \to v \to t$ e $s \to v \to u \to t$\\Aumento $\delta=1 \implies 2 \times 10^6$ passos (BFS / Dinic: 1 fase / 2 passos)};
\end{scope}
\end{tikzpicture}
%
}
\caption{Contraexemplo patológico de Ford-Fulkerson com $M = 10^6$ e arco central unitário $c=1$: (a) rede com fluxo ótimo $|f^*| = 2M$; (b) busca ingênua por DFS alternando caminhos em zigue-zague com aumento $\delta = 1$, demandando $2 \times 10^6$ iterações. Em contraste, métodos baseados em BFS (Edmonds-Karp e Dinic) priorizam caminhos mais curtos ($s \to u \to t$ e $s \to v \to t$), saturando a rede em 2 passos ou em 1 única fase.}
\label{fig:ford_fulkerson_patologico}
\end{figure}

\paragraph{Algoritmo de Edmonds-Karp (1972).} O método seleciona a cada iteração o caminho aumentante de menor quantidade de arcos através de Busca em Largura (BFS) em $D_f$~\cite{edmonds1972theoretical}. Demonstra-se que a distância residual mínima $d(s, v)$ é monotonicamente não decrescente ao longo das iterações. Entre duas saturações consecutivas de um mesmo arco $(u, v)$, a distância da origem sofre um acréscimo de no mínimo duas unidades ($d'(s, u) \ge d(s, u) + 2$). Como a distância máxima em um caminho simples é $|V|-1$, cada arco é saturado no máximo $|V|/2$ vezes. Sendo $|A|$ o número de arcos, ocorrem no máximo $\mathcal{O}(|V| \cdot |A|)$ aumentos. Com custo $\mathcal{O}(|A|)$ por BFS, a complexidade é $\mathcal{O}(|V| \cdot |A|^2)$~\cite{korte2018combinatorial}.

\paragraph{Algoritmo de Dinic (1970).} Yefim Dinitz estruturou o processamento em fases~\cite{dinic1970algorithm}. No início de cada fase, constrói-se o \textbf{Digrafo de Níveis} através de uma BFS global a partir de $s$, retendo apenas os arcos residuais que avançam de nível ($d(v) = d(u) + 1$). Em seguida, emprega-se uma DFS para saturar múltiplos caminhos em um único \textbf{Fluxo Bloqueador} (\textit{Blocking Flow}). Utilizando um vetor de ponteiros mortos (\textit{dead-end pointers}) que evita reexplorar arcos sem capacidade, a fase completa opera em $\mathcal{O}(|V| \cdot |A|)$. Como $d(s, t)$ cresce estritamente a cada fase, ocorrem no máximo $|V|-1$ fases, resultando em complexidade $\mathcal{O}(|V|^2 \cdot |A|)$. Em redes unitárias (como em grafos bipartidos), o número de fases limita-se a $\mathcal{O}(\sqrt{|V|})$, atingindo $\mathcal{O}(|A|\sqrt{|V|})$~\cite{ahuja1993network}.

\begin{figure}[!ht]
\centering
\begin{tikzpicture}[thick]
\foreach \x in {0,2.5,5.0,7.5} { \draw[layer line] (\x,-0.7) -- (\x,3.8); }

\node[gray,font=\footnotesize] at (0,4.1) {nív. 0};
\node[gray,font=\footnotesize] at (2.5,4.1) {nív. 1};
\node[gray,font=\footnotesize] at (5.0,4.1) {nív. 2};
\node[gray,font=\footnotesize] at (7.5,4.1) {nív. 3};

\node[source node] (s) at (0,1.5) {$s$};
\node[flow node] (a) at (2.5,3) {$a$};
\node[flow node] (b) at (2.5,0) {$b$};
\node[flow node] (c) at (5.0,3) {$c$};
\node[flow node] (d) at (5.0,0) {$d$};
\node[sink node] (t) at (7.5,1.5) {$t$};

\draw[flow edge,blue,line width=1.5pt] (s) -- (a);
\draw[flow edge,blue,line width=1.5pt] (s) -- (b);
\draw[flow edge,blue,line width=1.5pt] (a) -- (c);
\draw[flow edge,blue,line width=1.5pt] (b) -- (d);
\draw[flow edge,blue,line width=1.5pt] (a) -- (d);
\draw[flow edge,blue,line width=1.5pt] (b) -- (c);
\draw[flow edge,blue,line width=1.5pt] (c) -- (t);
\draw[flow edge,blue,line width=1.5pt] (d) -- (t);

\draw[nontree edge,bend right=30] (a) to (b);
\draw[nontree edge,bend right=20] (c) to (a);
\end{tikzpicture}

\caption{Digrafo de Níveis no algoritmo de Dinic: arcos azuis são admissíveis ($d(v) = d(u)+1$) e integram a fase de fluxo bloqueador; arcos cinza tracejados não avançam nível e são desconsiderados na fase.}
\label{fig:grafo_niveis}
\end{figure}

\subsection{Algoritmos de Pré-fluxo e Alturas (Push-Relabel)}

A família de algoritmos de Goldberg e Tarjan~\cite{goldberg1988new} adota uma perspectiva dual: relaxa-se a restrição de conservação de fluxo durante o cálculo, mantendo a ausência de caminhos aumentantes como invariante contínua. Um pré-fluxo $f$ respeita as capacidades ($0 \le f(a) \le c(a)$) e admite excessos nodais não negativos $e(v) \ge 0$ para todo $v \in V \setminus \{s, t\}$:
\begin{equation}
    e(v) = \sum_{\{u \mid (u, v) \in A\}} f(u, v) - \sum_{\{w \mid (v, w) \in A\}} f(v, w) \ge 0
\end{equation}

Um vértice com $e(v) > 0$ é denominado \textbf{ativo}. Para direcionar os excessos até o sumidouro $t$ (ou devolvê-los à fonte $s$), mantém-se uma função de alturas $h: V \to \mathbb{N}$ válida, que satisfaz: (i) $h(t) = 0$; (ii) $h(s) = |V|$; e (iii) $h(u) \le h(v) + 1$ para todo $(u, v) \in A_f$. A terceira condição impede a existência de caminhos de $s$ para $t$ em $D_f$. O algoritmo executa duas operações elementares sobre vértices ativos:
\begin{itemize}
    \item \textbf{Push (Empurrar):} Aplicado quando existe arco residual admissível ($c_f(u, v) > 0$ e $h(u) = h(v) + 1$). Envia-se $\delta = \min(e(u), c_f(u, v))$ unidades. Se $\delta = c_f(u, v)$, o push é saturante; caso contrário, é não saturante, zerando $e(u)$;
    \item \textbf{Relabel (Rerrotular):} Aplicado quando $u$ está ativo e não possui vizinhos residuais em nível inferior ($c_f(u, v) > 0 \implies h(u) \le h(v)$). Atualiza-se a altura para $h(u) = 1 + \min \{h(v) \mid c_f(u, v) > 0\}$.
\end{itemize}

Com política de seleção de vértices ativos em fila FIFO, a complexidade assintótica é $\mathcal{O}(|V|^3)$~\cite{goldberg1988new}.

\paragraph{Heurística de Lacuna (\textit{Gap Heuristic}).} Mantendo um vetor de contagem do número de vértices em cada estrato de altura $count[h]$, detecta-se se algum nível $k < |V|$ esvaziou-se ($count[k] = 0$). Como a condição $h(u) \le h(v) + 1$ limita a inclinação dos arcos residuais, nenhum vértice com $h(u) > k$ consegue alcançar o sumidouro ($h(t) = 0$). O algoritmo eleva suas alturas para $|V| + 1$, eliminando operações redundantes de relabel e devolvendo o fluxo diretamente à fonte. Associada à seleção por Maior Altura (\textit{Highest Label}), atinge-se a complexidade $\mathcal{O}(|V|^2 \sqrt{|A|})$~\cite{goldberg1988new}.

\begin{figure}[!ht]
\centering
\begin{tikzpicture}[thick]
  \draw[layer line] (-2, 0) -- (6.5, 0) node[right,black,font=\footnotesize] {$h=0$ (sumidouro $t$)};
  \foreach \y/\h in {1.3/1, 2.6/2, 4.8/4} {
    \draw[layer line] (-2, \y) -- (6.5, \y) node[right,black,font=\footnotesize] {$h=\h$};
  }
  \draw[red!50,dashed,line width=1.5pt] (-2, 3.7) -- (6.5, 3.7) node[right,red!80!black,font=\footnotesize] {$h=3$ (\textbf{LACUNA})};

  \node[sink node,font=\footnotesize] (t) at (2, 0) {$t$};
  \node[flow node,fill=yellow!20,font=\footnotesize] (a) at (2, 1.3) {$a$};
  \node[flow node,fill=yellow!20,font=\footnotesize] (b) at (0.5, 2.6) {$b$};
  \node[flow node,fill=yellow!20,font=\footnotesize] (c) at (3.5, 2.6) {$c$};

  \node[gap node,fill=red!20,font=\footnotesize] (u) at (2, 4.8) {$u$};

  \draw[->,blue!80!black] (b) -- (a);
  \draw[->,blue!80!black] (c) -- (a);
  \draw[->,blue!80!black] (a) -- (t);

  \draw[->,red,line width=1.4pt] (u) -- node[left,pos=0.5,font=\scriptsize,text=black,align=center,xshift=-5mm,fill=white,rounded corners=2pt,inner sep=2pt] {Queda de 2 níveis\\(inadmissível em $A_f$)} (b);
  \draw[->,red,line width=1.4pt] (u) -- node[right,pos=0.5,font=\scriptsize,text=black,align=center,xshift=5mm,fill=white,rounded corners=2pt,inner sep=2pt] {$h(u) \not\le h(c) + 1$\\(inadmissível em $A_f$)} (c);
\end{tikzpicture}

\caption{Mecanismo da Heurística de Lacuna: a restrição $h(u) \le h(v) + 1$ impede arcos residuais com salto de altura superior a uma unidade. O esvaziamento do nível intermediário $h=3$ comprova que vértices com $h \ge 4$ estão desconectados do sumidouro.}
\label{fig:gap_prova}
\end{figure}

\subsection{Reduções Combinatórias em Grafos Bipartidos}

A formulação de fluxo máximo permite resolver problemas clássicos em grafos bipartidos $G = (X \cup Y, E)$ através de reduções polinomiais exatas:

\begin{enumerate}[leftmargin=1.5em, itemsep=0.15em, topsep=0.2em, parsep=0em]
    \item \textbf{Emparelhamento Máximo Bipartido (\textit{Maximum Bipartite Matching}):} Constrói-se uma rede direcionada adicionando uma superfonte $s$ conectada a cada $x \in X$ com capacidade 1, e cada $y \in Y$ conectado ao supersumidouro $t$ com capacidade 1. Cada aresta $\{x, y\} \in E$ torna-se um arco $(x, y)$ com capacidade infinita (ou unitária). Pelo Teorema da Integralidade, o fluxo máximo tem valor inteiro e coincide rigorosamente com a cardinalidade máxima de emparelhamento: $|f^*| = \nu(G)$.

\begin{figure}[!ht]
\centering
\begin{tikzpicture}[thick]
\begin{scope}[xshift=0cm]
  \tikzsubcaption{(1.5,3.8)}{(a) Grafo bipartido original}
  \foreach \i/\y in {1/3, 2/1.5, 3/0} {
    \node[bip L] (l\i) at (0,\y) {$l_\i$};
    \node[bip R] (r\i) at (3,\y) {$r_\i$};
  }
  \draw[-] (l1) -- (r1); \draw[-] (l1) -- (r2);
  \draw[-] (l2) -- (r2); \draw[-] (l2) -- (r3);
  \draw[-] (l3) -- (r1); \draw[-] (l3) -- (r3);
\end{scope}
\node at (4.8,1.5) {\Huge $\Rightarrow$};

\begin{scope}[xshift=6.2cm]
  \tikzsubcaption{(2.5,3.8)}{(b) Rede de fluxo equivalente $D'$}
  \node[source node] (s) at (0,1.5) {$s$};
  \node[sink node] (t) at (5,1.5) {$t$};
  \foreach \i/\y in {1/3, 2/1.5, 3/0} {
    \node[bip L,font=\footnotesize] (fl\i) at (1.8,\y) {$l_\i$};
    \node[bip R,font=\footnotesize] (fr\i) at (3.2,\y) {$r_\i$};
    \draw[flow edge] (s) -- node[above,font=\footnotesize] {$1$} (fl\i);
    \draw[flow edge] (fr\i) -- node[above,font=\footnotesize] {$1$} (t);
  }
  \draw[flow edge,blue,line width=1.5pt] (fl1) -- (fr1); \draw[flow edge] (fl1) -- (fr2);
  \draw[flow edge,blue,line width=1.5pt] (fl2) -- (fr2); \draw[flow edge] (fl2) -- (fr3);
  \draw[flow edge] (fl3) -- (fr1); \draw[flow edge,blue,line width=1.5pt] (fl3) -- (fr3);
\end{scope}
\end{tikzpicture}

\caption{Redução do problema de Emparelhamento Máximo Bipartido para Fluxo Máximo: (a) grafo bipartido original; (b) rede com superfonte $s$, supersumidouro $t$ e capacidades unitárias. Os arcos azuis com $f(u, v) = 1$ indicam as arestas do emparelhamento máximo $|f^*| = 3$.}
\label{fig:fluxo_bipartido}
\end{figure}

    \item \textbf{Teorema de König (1931) e Cobertura Mínima de Vértices:} Em todo grafo bipartido, a cardinalidade do emparelhamento máximo é igual à cardinalidade da cobertura mínima de vértices ($\tau(G)$):
    \begin{equation}
        \nu(G) = \tau(G)
    \end{equation}
    A demonstração constrói o corte mínimo residual $(S^*, T^*)$ em $D_{f^*}$ a partir de $s$. O conjunto $C^* = (X \setminus S^*) \cup (Y \cap S^*)$ cobre todas as arestas do grafo e sua cardinalidade satisfaz $|C^*| = c(S^*, T^*) = |f^*| = \nu(G)$, fornecendo um algoritmo polinomial em tempo $\mathcal{O}(|E|\sqrt{|V|})$ via Dinic.

    \item \textbf{Conjunto Independente Máximo e Identidades de Gallai (1959):} Um subconjunto de vértices $I$ é independente se e somente se seu complementar $V \setminus I$ for uma cobertura de vértices. Pelo Teorema de Gallai, $\alpha(G) + \tau(G) = |V|$ e $\nu(G) + \rho(G) = |V|$, onde $\alpha(G)$ é a cardinalidade do conjunto independente máximo e $\rho(G)$ da cobertura mínima de arestas. Em grafos bipartidos, $\alpha(G) = \rho(G) = |V| - \nu(G)$, e o conjunto ótimo é extraído diretamente por $I^* = V \setminus C^*$ em tempo $\mathcal{O}(|V|)$.

    \item \textbf{Teorema do Casamento de Hall (1935):} Um grafo bipartido $G = (X \cup Y, E)$ possui um emparelhamento que cobre todos os vértices de $X$ se e somente se $|N(S)| \ge |S|$ para todo $S \subseteq X$, onde $N(S)$ é a vizinhança de $S$. A prova deduz-se diretamente de König aplicando a contradição de corte mínimo.
\end{enumerate}

\subsection{Conectividade e Roteamento}

A determinação do número de rotas independentes entre dois vértices $s$ e $t$ reduz-se a fluxo máximo:
\begin{itemize}[leftmargin=1.5em, itemsep=0.15em, topsep=0.2em, parsep=0em]
    \item \textbf{Caminhos Disjuntos por Arcos:} Fixando $c(a) = 1$ para todo arco, o valor $|f^*|$ equivale ao número máximo de rotas direcionadas sem arcos em comum (Teorema de Menger por arcos, 1927).
    \item \textbf{Caminhos Disjuntos por Vértices e \textit{Vertex Splitting}:} Para impedir que rotas compartilhem vértices intermediários, divide-se cada vértice $v \in V \setminus \{s, t\}$ em dois nós, $v_{\text{in}}$ e $v_{\text{out}}$, conectados por um arco interno $(v_{\text{in}}, v_{\text{out}})$ com capacidade unitária $c(v_{\text{in}}, v_{\text{out}}) = 1$. O corte mínimo nesta rede expandida identifica o conjunto de vértices cuja remoção desconecta a fonte do sumidouro (Teorema de Menger por vértices, 1927).
\end{itemize}

\subsection{Teoria do Fluxo de Custo Mínimo (MCF)}

No problema de Fluxo de Custo Mínimo, associa-se a cada arco $(u, v) \in A$ uma capacidade $c(u, v) \ge 0$ e um custo linear unitário $w(u, v) \in \mathbb{R}$. Cada vértice $v \in V$ possui demanda $b(v) \in \mathbb{R}$ com $\sum_{v \in V} b(v) = 0$. O problema formula-se pelo programa linear:
\begin{align}
    \min_{f} \quad & \sum_{(u, v) \in A} w(u, v) \cdot f(u, v) \\
    \sa \quad & \sum_{\{w \mid (v, w) \in A\}} f(v, w) - \sum_{\{u \mid (u, v) \in A\}} f(u, v) = b(v), \quad \forall v \in V \\
    & 0 \le f(u, v) \le c(u, v), \quad \forall (u, v) \in A
\end{align}

Na rede residual $D_f$, os arcos diretos preservam o custo $w_f(u, v) = w(u, v)$ e os arcos reversos recebem custo antissimétrico $w_f(v, u) = -w(u, v)$, refletindo a devolução de custo ao cancelar fluxo.

\begin{prop}[Condição de Otimalidade por Ciclos Negativos --- Klein, 1967]\label{prop:klein_rel}
Um fluxo viável $f$ é de custo mínimo se, e somente se, o digrafo residual associado $D_f$ não contém nenhum ciclo direcionado de custo total estritamente negativo.
\end{prop}

\begin{prop}[Custos Reduzidos e Potenciais Nodais]\label{prop:potenciais_rel}
Seja $\pi: V \to \mathbb{R}$ um vetor de potenciais nodais (variáveis duais das restrições de balanço). O custo reduzido de um arco residual $(u, v) \in A_f$ é definido por:
\begin{equation}
    \bar{w}_\pi(u, v) = w_f(u, v) - \pi(u) + \pi(v)
\end{equation}
Um fluxo viável $f$ é ótimo se e somente se existe $\pi$ tal que $\bar{w}_\pi(u, v) \ge 0$ para todo $(u, v) \in A_f$.
\end{prop}

A não negatividade dos custos reduzidos corresponde às condições de Folga Complementar de Karush-Kuhn-Tucker (KKT). Mantendo os potenciais atualizados ao longo das iterações, viabiliza-se o uso do algoritmo de Dijkstra com fila de prioridades no lugar de algoritmos com suporte a custos negativos no método \textit{Successive Shortest Path} (SSP).

\subsection{O Algoritmo Network Simplex}

O algoritmo \textit{Network Simplex} especializa o método Simplex linear explorando a estrutura combinatória de árvores geradoras da rede (veja~\cite{ahuja1993network, dantzig1951application}). Uma solução básica viável particiona os arcos em três classes $(T, L, U)$:
\begin{enumerate}
    \item $T$: arcos básicos que formam uma árvore geradora do digrafo;
    \item $L$: arcos não básicos com fluxo no limite inferior ($f(a) = 0$);
    \item $U$: arcos não básicos com fluxo saturado no limite superior ($f(a) = c(a)$).
\end{enumerate}

\begin{prop}[Total Unimodularidade (TUM) da Matriz Nó-Arco]\label{prop:tum_rel}
Seja $\mathbf{M} \in \{-1, 0, 1\}^{|V| \times |A|}$ a matriz de incidência nó-arco de $D$. Todo determinante de submatriz quadrada de $\mathbf{M}$ pertence a $\{-1, 0, 1\}$. Consequentemente, para qualquer vetor de demandas e capacidades inteiras, toda solução básica (base em árvore) é puramente inteira.
\end{prop}

A Total Unimodularidade garante que a matriz básica $\mathbf{B}$ associada a qualquer árvore $T$ satisfaça $|\det \mathbf{B}| = 1$, eliminando divisões na Regra de Cramer e garantindo aritmética inteira exata.

Cada iteração do Network Simplex executa quatro passos: (i) \textbf{Pricing:} identificação de arco violador $(u, v)$ fora da árvore ($\bar{w}_\pi < 0$ em $L$ ou $\bar{w}_\pi > 0$ em $U$); (ii) \textbf{Ciclo Fundamental:} determinação do ciclo $C$ formado por $(u, v)$ em $T$ via busca do ancestral comum mais próximo (\textit{LCA}); (iii) \textbf{Gargalo:} cômputo da capacidade residual mínima $\delta$ em $C$ e determinação do arco que deixa a base; e (iv) \textbf{Pivoteamento:} atualização do fluxo e ajuste dos potenciais nodais na subárvore em tempo $\mathcal{O}(|V|)$ com auxílio de listas de pré-ordem (\lstinline{thread}).

\begin{figure}[!ht]
\centering
\begin{tikzpicture}[thick]
\begin{scope}[xshift=0cm]
  \tikzsubcaption{(2.0,3.6)}{(a) Entrada do arco $(3,4)$ no ciclo fundamental}
  \node[source node] (a1) at (0, 1.3) {$1$};
  \node[flow node] (a2) at (2, 2.6) {$2$};
  \node[flow node] (a3) at (2, 0.0) {$3$};
  \node[sink node] (a4) at (4, 1.3) {$4$};

  \draw[tree edge] (a1) -- (a2);
  \draw[tree edge] (a2) -- (a4);
  \draw[tree edge] (a1) -- (a3);
  \draw[pivot enter] (a3) -- node[below=3pt,sloped,font=\scriptsize,text=red!80!black] {entra} (a4);
\end{scope}

\node at (5.2,1.3) {\Large $\Rightarrow$};

\begin{scope}[xshift=6.4cm]
  \tikzsubcaption{(2.0,3.6)}{(b) Nova base em árvore após o pivot}
  \node[source node] (b1) at (0, 1.3) {$1$};
  \node[flow node] (b2) at (2, 2.6) {$2$};
  \node[flow node] (b3) at (2, 0.0) {$3$};
  \node[sink node] (b4) at (4, 1.3) {$4$};

  \draw[tree edge] (b1) -- (b2);
  \draw[pivot leave] (b2) -- node[above=3pt,sloped,font=\scriptsize,text=gray!80!black] {sai} (b4);
  \draw[tree edge] (b1) -- (b3);
  \draw[tree edge] (b3) -- (b4);
\end{scope}
\end{tikzpicture}

\caption{Iteração do Network Simplex: (a) o arco não básico violador $(3,4)$ com $\bar{w}_\pi < 0$ fecha um ciclo fundamental em $T$; (b) o fluxo é aumentado por $\delta$ e o arco saturado $(2,4)$ deixa a base, restabelecendo a base em árvore.}
\label{fig:simplex_pivoteamento}
\end{figure}
